\section{The primitive genus-three count}\label{sec:pairs}

Put $D=3s+4f$ on $F\simeq\FF_1$. The intersection numbers are
\[
 s^2=-1,\quad sf=1,\quad f^2=0,\quad K_F=-2s-3f,
 \qquad D^2=15,\quad K_FD=-11.
\]
Thus $p_a(D)=3$. In the presentation $F=\Bl_p\PP^2$,
$D=4H-E_p$; consequently
\begin{equation}\label{eq:linear-system}
 h^0(F,\mathcal O_F(D))=5+4+3+2=14,\qquad |D|\simeq\PP^{13}.
\end{equation}
We must control all stable pairs in class $\alpha$, including
curves with nonreduced structure transverse to $F$.

\subsection{A surface inequality}
For $a,b\ge0$, the Euler characteristic of a divisor of class
$as+bf$ is
\begin{equation}\label{eq:phi}
 \phi(a,b)=\frac{a(a+1)}2+b(1-a),
 \qquad \phi(3,4)=-2.
\end{equation}
This follows from the surface Riemann--Roch formula
$\chi(\mathcal O_C)=-C(C+K_F)/2$.

\begin{lemma}\label{lem:split}
If $\sum_i(a_i,b_i)=(3,4)$ is a sum of at least two nonzero
pairs of nonnegative integers, then $\sum_i\phi(a_i,b_i)\ge0$.
\end{lemma}
\begin{proof}
Partition the positive $a_i$'s. If the partition is $(3)$,
let $t\ge1$ be the total $b$-coefficient on the other terms.
The sum is $-2+3t\ge1$. For the partition $(2,1)$, write
$b_A,b_B$ for the two $b$-coefficients and $t$ for the remaining
total. The sum is $4-b_A+t=b_B+2t\ge0$.
For $(1,1,1)$ the sum is $3+t\ge3$. These are all possibilities.
\end{proof}

\subsection{Transverse thickenings}
\begin{proposition}\label{prop:minimum}
Every pure one-dimensional Cohen--Macaulay subscheme $Z\subset Y$
of class $\alpha$ satisfies $\chi(\mathcal O_Z)\ge-2$.
Equality forces $Z$ to be a divisor in $|D|$ on $F$, with its
scheme structure contained in $F$.
\end{proposition}
\begin{proof}
The cycle of $Z$ is supported on $F$ by Proposition~\ref{prop:support}.
The filtration by powers of its ideal $I_F$ is finite. Since
$N_{F/Y}=K_F$, its graded quotients have the form
\begin{equation}\label{eq:layers}
 G_j=I_F^j\mathcal O_Z/I_F^{j+1}\mathcal O_Z
 \simeq K_F^{-j}\otimes(\mathcal O_F/J_j)
\end{equation}
for ideal sheaves $J_j\subset\mathcal O_F$. On the smooth surface,
the divisorial hull of $J_j$ is $\mathcal O_F(-D_j)$ for an
effective divisor $D_j$. The quotient
$\mathcal O_F(-D_j)/J_j$ has finite length $\ell_j\ge0$.
Additivity of Euler characteristic gives
\begin{equation}\label{eq:layer-chi}
 \chi(G_j)=\chi(\mathcal O_{D_j})-jK_FD_j+\ell_j,
 \qquad\sum_j[D_j]=[D].
\end{equation}
The last equality holds in the divisor-class group of $F$, since
its two coefficients are detected by $H_2,H_3$. Each twist
$-jK_FD_j$ is nonnegative because $-K_F$ is ample.

If at least two $D_j$ are nonzero, Lemma~\ref{lem:split} shows
$\chi(\mathcal O_Z)\ge0$. Otherwise there is exactly one
nonzero divisorial layer. At the generic point of every component,
a nonzero higher layer requires a nonzero zeroth layer, so the
unique divisorial layer must be the zeroth one. Then
$I_F\mathcal O_Z$ is supported at finitely many points.
As a subsheaf of the pure sheaf $\mathcal O_Z$, it is zero.
Thus $Z$ is contained in $F$ scheme-theoretically. A pure curve
on a smooth surface is an effective Cartier divisor, so
$Z\in|D|$ and $\chi(\mathcal O_Z)=-2$.
This also excludes a transverse thickening confined to finitely
many points, which cycle support alone would not detect.
\end{proof}

For a stable pair $\mathcal O_Y\to\mathcal F$, its image is
$\mathcal O_Z$ for a pure curve and its cokernel is
zero-dimensional. The exact sequence
\[
 0\longrightarrow\mathcal O_Z\longrightarrow\mathcal F
 \longrightarrow Q\longrightarrow0
\]
therefore gives $\chi(\mathcal F)\ge-2$. At equality,
$Q=0$ and $Z\in|D|$.

\subsection{The minimal moduli scheme}
\begin{proposition}\label{prop:pt-count}
The stable-pair moduli space $P_{-2}(Y,\alpha)$ is
scheme-theoretically $\PP^{13}$. Its virtual count is $-14$, and
$P_n(Y,\alpha)$ is empty for $n<-2$.
\end{proposition}
\begin{proof}
The set-theoretic assertions follow from Proposition~\ref{prop:minimum}.
Every pair section at Euler characteristic $-2$ is surjective.
Fiberwise surjectivity and Nakayama's lemma identify the pair
functor with the corresponding Hilbert functor. The surface divisors
give a closed immersion $|D|\hookrightarrow P_{-2}(Y,\alpha)$
which exhausts the closed points.

For every $C\in|D|$, including nonreduced divisors, the normal
sequence for the regular embeddings $C\subset F\subset Y$ is
\begin{equation}\label{eq:normal}
 0\longrightarrow\mathcal O_C(D)\longrightarrow N_{C/Y}
 \longrightarrow K_F|_C\longrightarrow0.
\end{equation}
We have $H^0(F,K_F)=0$ and
\[
 H^1(F,K_F-D)=H^1(F,D)^*=0,
\]
by Serre duality and Kodaira vanishing, since
$D-K_F=5s+7f$ is ample. Restriction of $K_F$ to $C$ therefore
has no sections. Rationality of $F$ and \eqref{eq:linear-system}
give $h^0(\mathcal O_C(D))=13$. By \eqref{eq:normal} the
ambient Hilbert tangent space has dimension thirteen at every
point of the closed $\PP^{13}$. Its local dimension is at least
thirteen, so its local rings are regular of dimension thirteen.
The closed immersion is consequently an isomorphism.

The Calabi--Yau stable-pair obstruction theory is symmetric. On a
smooth proper moduli space of dimension $r$ its virtual degree is
$(-1)^r\chi_{\mathrm{top}}$ \cite[p.~3]{PT}. Hence
\begin{equation}\label{eq:minus-fourteen}
 P_{-2,\alpha}=(-1)^{13}\chi_{\mathrm{top}}(\PP^{13})=-14.
\end{equation}
\end{proof}

\subsection{Passing to the connected coefficient}
Taking the logarithm of the geometric PT series adds products
over at least two nonzero effective summands of $\alpha$ to its
coefficient at $\alpha$. Each constituent cycle lies in $F$.
Collect the divisorial layers \eqref{eq:layers} of all constituent
curves. They give at least two nonzero effective terms with total
class $D$. Lemma~\ref{lem:split} and \eqref{eq:layer-chi} show
that the sum of their Euler characteristics is nonnegative.
The twists, zero-dimensional torsion and pair cokernels only
increase this sum. Thus none of these lower products contributes
a negative power of the PT variable $z$.

It follows that the connected coefficient has no term below
$z^{-2}$ and has coefficient $-14$ at $z^{-2}$. By primitivity,
Lemma~\ref{lem:factor} and \eqref{eq:pt-one} give exactly
\begin{equation}\label{eq:primitive-pt}
 [Q^\alpha]\log\PT(\Z_Y)
 =\sum_{g\ge0}e_{\alpha,g}z^{1-g}(1+z)^{2g-2}.
\end{equation}
This sum has finite genus support. Its largest nonzero genus
contributes a unique lowest Laurent power, with coefficient one.
Consequently
\begin{equation}\label{eq:primitive-count}
 e_{\alpha,g}=0\ (g>3),\qquad
 e_{\alpha,3}=c_{\alpha,3}=-14.
\end{equation}
